% Generated draft: requires mathematical and source review.
\documentclass[11pt]{article}
\usepackage[margin=25mm]{geometry}
\usepackage{fontspec}
\setmainfont{Latin Modern Roman}
\usepackage{amsmath,amssymb,mathtools,unicode-math}
\setmathfont{Latin Modern Math}
\usepackage{enumitem,needspace,xurl,hyperref}
\setlength{\parindent}{0pt}
\setlength{\parskip}{5pt}
\setlength{\emergencystretch}{3em}
\newcommand{\N}{\mathbb{N}}
\newcommand{\Z}{\mathbb{Z}}
\newcommand{\Q}{\mathbb{Q}}
\newcommand{\R}{\mathbb{R}}
\newcommand{\C}{\mathbb{C}}
\newcommand{\F}{\mathbb{F}}
\newcommand{\A}{\mathbb{A}}
\newcommand{\PP}{\mathbb{P}}
\newcommand{\E}{\mathbb{E}}
\newcommand{\eps}{\varepsilon}
\newcommand{\dd}{\,\mathrm{d}}
\begin{document}
\section*{UnsolvedMath ID: 30001708}
\textbf{Source problem code:} catalogue entry
\subsection{Definitions and mathematical statement}
For notation only, let \(G\) denote the unspecified Lie group underlying the source question, let \(\mathfrak g\) be its Lie algebra, and let \(\mathfrak g^*\) be its dual. A coadjoint orbit through \(\xi\in\mathfrak g^*\) is
\[
O_\xi=\{\operatorname{Ad}^*(g)\xi:g\in G\}.
\]
The source does not specify the group, the ground field, or hypotheses under which coadjoint elements are classified as nilpotent or semisimple. Those terms must therefore be interpreted in an appropriate, still unspecified setting. Write \(Q(O)\) for the representation assigned to an orbit \(O\) by the intended geometric quantization procedure, whenever that procedure is defined; the source does not specify that procedure or its auxiliary choices.

\textbf{Construction problem.} Given a nilpotent coadjoint orbit \(O_{\mathrm{nilp}}\) approached by a family of semisimple coadjoint orbits \((O_\nu)_{\nu\in I}\), and given their geometric quantizations \(Q(O_\nu)\), construct the representation
\[
Q(O_{\mathrm{nilp}})
\quad\text{from the family}\quad
\bigl(Q(O_\nu)\bigr)_{\nu\in I}.
\]
Here \(I\) merely denotes an unspecified parameter set. The meaning of ``approached,'' the limiting parameter, the admissibility conditions for quantization, and the operation that constructs the target representation are not supplied. In particular, this formulation does not assert an equality with a limit of representations in any specified topology. The supplied assessment describes the intended scope as arbitrary nilpotent orbits, but the precise quantifiers over groups and approximating families are unresolved. These omissions prevent a fully specified mathematical assertion without further source clarification.
\subsection{Short English statement}
The problem asks how to construct a representation quantizing a nilpotent coadjoint orbit from representations obtained by quantizing semisimple orbits that approach it. The ambient group, the meaning of approach, and the quantization and reconstruction procedures remain unspecified in the supplied record.
\subsection{Sources}
\begin{itemize}
\item Local registry: \texttt{data/unsolvedmath/problems.json}, permanent UnsolvedMath ID \(30001708\), title ``Quantization of Nilpotent Coadjoint Orbits by Degeneration,'' exact source problem code \texttt{catalogue entry}. The supplied record belongs to the Oberwolfach Reports Open Problems collection.
\item Supplied report citation: \emph{Geometric Quantization in the Non-compact Setting} (2011), Oberwolfach Reports, pp.\ 425--521, DOI \texttt{10.4171/owr/2011/09}. Source URL and supplied literature-source link, labelled ``original Oberwolfach report'': \url{https://doi.org/10.4171/owr/2011/09}. The background describes these reports as records of mathematical workshops and their open-problem discussions. No individual proposer is identified; the supplied proposed year is 2011.
\item Supplied literature assessment dated \texttt{2026-08-21}: the record labels the problem open and states that quantization for arbitrary nilpotent orbits remains a problem, with constructions available only in selected minimal-orbit cases. This is supplied dated curation, explicitly requiring independent verification of time-sensitive status and cited claims; it has not been independently verified here. The only listed literature source is the report linked above, and no particular minimal-orbit construction is cited.
\end{itemize}
\end{document}
